\documentclass{ctexart}
\usepackage{template}
\usepackage{esint}
\begin{document}\pagestyle{empty}
\begin{center}\Large
北京大学2025-26学年第1学期高数B期中试题
\end{center}

\begin{enumerate}[leftmargin=*,label=\textbf{\arabic*.},topsep=0pt,parsep=0pt,itemsep=0pt,partopsep=0pt]
    \item \textbf{(10\songti{分})}\ 求序列极限.
\begin{enumerate}[label=\tbf{(\alph*)},leftmargin=*,topsep=0pt,parsep=0pt,itemsep=0pt,partopsep=0pt]
\item \[\lim_{n\to\infty}\sqrt[n]{3^{2n}+4^n}\]
\item \[\lim_{n\to\infty}\sum_{k=1}^n\dfrac{(n+k)^2}{n^3}\]
\end{enumerate}
    \item \textbf{(10\songti{分})}\ 求函数极限.
\begin{enumerate}[label=\tbf{(\alph*)},leftmargin=*,topsep=0pt,parsep=0pt,itemsep=0pt,partopsep=0pt]
\item \[\lim_{x\to4}\dfrac{\sqrt{2x+1}-3}{\sqrt{x-2}-\sqrt2}\]
\item \[\lim_{x\to+\infty}\left(\sin\sqrt{x+1}-\sin\sqrt x\right)\]
\end{enumerate}
    \item \textbf{(10\songti{分})}\ 求导数.
\begin{enumerate}[label=\tbf{(\alph*)},leftmargin=*,topsep=0pt,parsep=0pt,itemsep=0pt,partopsep=0pt]
\item 求$\left(x^{x^a}\right)'$, $x>0$,其中$a>0$是常数.
\item 求$\displaystyle\dfrac{\di}{\di x}\int_{\sin x}^{\tan x}\sqrt{1+t^2}\di t$.
\end{enumerate}
    \item \textbf{(10\songti{分})}\ 求多阶导数.
\begin{enumerate}[label=\tbf{(\alph*)},leftmargin=*,topsep=0pt,parsep=0pt,itemsep=0pt,partopsep=0pt]
\item 设函数$y=y(x)$由参数方程
\[\begin{cases}x=R(t-\sin t),\\ y=R(1-\cos t),\end{cases}\qquad t\in(0,2\pi)\]
确定,其中$R>0$是常数.求$y'(x)$, $y''(x)$.
\item 设$f(x)=x^2\e^{-x^2}$.求$f^{(n)}(x)$, $n=1,2,3$.
\end{enumerate}
    \item \textbf{(18\songti{分})}\ 求不定积分.
\begin{enumerate}[label=\tbf{(\alph*)},leftmargin=*,topsep=0pt,parsep=0pt,itemsep=0pt,partopsep=0pt]
\item \[\int x\ln\sqrt{1+x^2}\di x\]
\item \[\int\dfrac{x^2+2x+3}{(x+2)^2(x+3)^2}\di x\]
\item \[\int x^2\sqrt{1-x^2}\di x\]
\end{enumerate}
    \item \textbf{(12\songti{分})}\ 求定积分.
\begin{enumerate}[label=\tbf{(\alph*)},leftmargin=*,topsep=0pt,parsep=0pt,itemsep=0pt,partopsep=0pt]
\item \[\int_{1/2}^1\sqrt{x(1-x)}\di x\]
\item \[\int_0^{\pi/2}\dfrac{\sin^2x}{\sin x+\cos x}\di x\]
\end{enumerate}
    \item \textbf{(20\songti{分})}\ 设$D$是由抛物线$y=2x^2$与直线$y=8$所围成的有界区域,$B$是由区域$D$绕$y$轴旋转一周所形成的旋转体.
\begin{enumerate}[label=\tbf{(\alph*)},leftmargin=*,topsep=0pt,parsep=0pt,itemsep=0pt,partopsep=0pt]
\item 求区域$D$的面积$A$.
\item 求旋转体$B$的体积$V$.
\item 求$B$的表面积$S$.
\end{enumerate}
    \item \textbf{(10\songti{分})}\ 设函数$f(x)=x^8-x^4-\cos x$, $x\in\R$.判断$f(x)$在$\R$上有多少个零点,给出证明.
\end{enumerate}
\end{document}
